% ============================================================
% Chapter 5: The 8 Iterations of Prompt Development
% ============================================================
\section{Iterative Prompt Development: 8 Runs to Final Validation}
\label{sec:iterations}

This section documents every iteration of the prompt development process in full detail. Each run was measured against the 200-row human gold standard. The complete history is also available in \texttt{reports/prompt\_version\_log.csv}.

\subsection{Summary of All 8 Iterations}

\begin{table}[H]
\centering
\caption{Complete 8-Iteration Prompt Development Log}
\label{tab:iterations}
\begin{tabular}{@{}clp{5.5cm}ccl@{}}
\toprule
\textbf{Run} & \textbf{Step} & \textbf{Key Change} & \textbf{Threat F1} & \textbf{Benefit F1} & \textbf{Gate} \\
\midrule
1 & 2.1 & Baseline v1 instructions & 0.643 & 0.286 & \failmark \\
2 & 2.1 & +Logistics/statistics exclusions & 0.652 & 0.857 & \warnmark \\
3 & 2.1 & +Strict NO rules + reminder & 0.438 & 0.667 & \failmark \\
4 & 2.1 & +International/non-immigration exclusions & 0.622 & 0.600 & \failmark \\
5 & 3.2 & Few-shot + CoT + encoding fix & 0.727 & 0.500 & \failmark \\
6 & 3.2 & +T5 pressure +B5 degree +B4 & 0.833 & 0.769 & \warnmark \\
7 & 3.2 & +Benefit NO logistical example & 0.857 & 0.769 & \warnmark \\
\rowcolor{pass!15} \textbf{8} & \textbf{3.2} & \textbf{+Benefit NO legal-status example} & \textbf{0.811} & \textbf{0.769} & \passmark \\
\bottomrule
\end{tabular}
\end{table}

% ------------
\subsection{Run 1: Baseline (FAILED)}
\label{sec:run1}

\textbf{Step:} 2.1 (Zero-Shot) \quad \textbf{Attempts to run:} 1

\textbf{Instructions:} Standard bacchuss routine with T1--T5 and B1--B5 sub-criteria. No exclusion rules, no few-shot examples.

\textbf{Results:}
\begin{itemize}
    \item \textbf{Threat:} Precision = 0.487, Recall = 0.947, F1 = \textbf{0.643}
    \item \textbf{Benefit:} Precision = 0.167, Recall = 1.000, F1 = \textbf{0.286}
\end{itemize}

\textbf{Failure Analysis:}
\begin{itemize}
    \item Both dimensions showed \textbf{perfect or near-perfect Recall} but \textbf{extremely low Precision}.
    \item The model was ``over-coding''---labelling far too many paragraphs as YES.
    \item For Threat, only 48.7\% of YES predictions were correct (Precision = 0.487).
    \item For Benefit, only 16.7\% of YES predictions were correct (Precision = 0.167).
    \item The model confused the \textit{presence of refugees} in a paragraph with economic framing.
\end{itemize}

\textbf{Action:} Add explicit exclusion rules to reduce false positives.

% ------------
\subsection{Run 2: Logistics \& Statistics Exclusions (PARTIAL PASS)}
\label{sec:run2}

\textbf{Step:} 2.1 (Zero-Shot) \quad \textbf{Attempts to run:} 1

\textbf{Changes:}
\begin{itemize}
    \item Added Threat exclusion: ``Logistics reporting---text only states that accommodation WAS set up with NO mention of cost = NO''
    \item Added Threat exclusion: ``Statistics reporting---bed counts, arrivals without burden language = NO''
    \item Added Benefit exclusion: ``Factual entitlement announcements---free transport stated without framing as benefit = NO''
\end{itemize}

\textbf{Results:}
\begin{itemize}
    \item \textbf{Threat:} Precision = 0.556, Recall = 0.789, F1 = \textbf{0.652}
    \item \textbf{Benefit:} Precision = 0.750, Recall = 1.000, F1 = \textbf{0.857}
\end{itemize}

\textbf{Analysis:}
\begin{itemize}
    \item Benefit F1 jumped dramatically from 0.286 to 0.857---above the 0.80 gate.
    \item Threat improved slightly but remained below gate (0.652).
    \item Benefit Precision tripled from 0.167 to 0.750.
\end{itemize}

\textbf{Gate:} Benefit = \passmark $|$ Threat = \failmark

% ------------
\subsection{Run 3: Strict NO Rules (REGRESSION --- FAILED)}
\label{sec:run3}

\textbf{Step:} 2.1 (Zero-Shot) \quad \textbf{Attempts to run:} 1

\textbf{Changes:} Added strict ``When in doubt, code NO'' rules and a tightened reminder for both dimensions.

\textbf{Results:}
\begin{itemize}
    \item \textbf{Threat:} Precision = 0.538, Recall = 0.368, F1 = \textbf{0.438}
    \item \textbf{Benefit:} Precision = 0.500, Recall = 1.000, F1 = \textbf{0.667}
\end{itemize}

\textbf{Critical Regression:}
\begin{itemize}
    \item \textcolor{fail}{\textbf{Threat F1 dropped from 0.652 to 0.438}} --- the worst score in the entire development process.
    \item Threat Recall collapsed from 0.789 to 0.368. The ``strict NO'' rule made the model too conservative, causing it to miss genuine Threat cases.
    \item Benefit also regressed from 0.857 to 0.667. The strictness was counterproductive.
\end{itemize}

\textbf{Lesson Learned:} Overly strict ``When in doubt, code NO'' instructions cause the model to under-code, destroying Recall. The instructions must balance Precision and Recall carefully.

\textbf{Action:} Reverted the strict NO rules. Adopted a more targeted approach: adding specific exclusion categories rather than blanket conservatism.

% ------------
\subsection{Run 4: International \& Non-Immigration Exclusions (FAILED)}
\label{sec:run4}

\textbf{Step:} 2.1 (Zero-Shot) \quad \textbf{Attempts to run:} 1

\textbf{Changes:}
\begin{itemize}
    \item Added exclusions for international politics and foreign aid (UK-Rwanda, EU-Hungary)
    \item Added exclusions for criminal activity (embezzlement, fraud)
    \item Added exclusions for non-immigration capacity limits (university spots, sports subsidies)
    \item Added exclusions for German domestic economic issues not linked to immigration
\end{itemize}

\textbf{Results:}
\begin{itemize}
    \item \textbf{Threat:} Precision = 0.538, Recall = 0.737, F1 = \textbf{0.622}
    \item \textbf{Benefit:} Precision = 0.429, Recall = 1.000, F1 = \textbf{0.600}
\end{itemize}

\textbf{Analysis:} Marginal improvement for Threat Recall (0.368 $\rightarrow$ 0.737) after reverting the strict rules, but F1 still below gate. Benefit continued to decline (0.667 $\rightarrow$ 0.600).

\textbf{Decision:} Zero-shot approach exhausted. Transition to Step~3.2 (few-shot + chain-of-thought reasoning).

% ------------
\subsection{Run 5: Few-Shot + CoT + Encoding Fix (FAILED)}
\label{sec:run5}

\textbf{Step:} 3.2 (Structured Inference) \quad \textbf{Attempts to run:} Multiple (encoding errors)

\textbf{Changes:}
\begin{itemize}
    \item Switched from zero-shot to few-shot with chain-of-thought (CoT) reasoning
    \item Added 5 training examples per dimension from \texttt{fewshot\_training\_blocks.csv}
    \item Implemented \texttt{clean\_text\_api()} to fix non-ASCII encoding errors
    \item Created custom \texttt{annotate\_direct()} function with 3-attempt retry logic per row
\end{itemize}

\textbf{Technical Failures During This Run:}
\begin{itemize}
    \item \textcolor{fail}{HTTP 400 errors} from the university API due to non-ASCII characters (umlauts, special quotes) in translated text
    \item Required development of the \texttt{clean\_text\_api()} helper
    \item Multiple script restarts before the encoding fix was implemented
\end{itemize}

\textbf{Results (after encoding fix):}
\begin{itemize}
    \item \textbf{Threat:} Precision = 0.857, Recall = 0.632, F1 = \textbf{0.727}
    \item \textbf{Benefit:} Precision = 1.000, Recall = 0.333, F1 = \textbf{0.500}
\end{itemize}

\textbf{Analysis:}
\begin{itemize}
    \item Threat Precision improved dramatically (0.538 $\rightarrow$ 0.857), but Recall dropped.
    \item Benefit Precision reached perfect 1.000 but Recall collapsed to 0.333---the model missed 2 out of 3 true Benefit cases.
    \item The few-shot examples were too restrictive for Benefit, causing under-coding.
\end{itemize}

% ------------
\subsection{Run 6: +T5 Financial Pressure, +B5 Degree, +B4 (ALMOST)}
\label{sec:run6}

\textbf{Step:} 3.2 \quad \textbf{Attempts to run:} 1

\textbf{Changes:}
\begin{itemize}
    \item Added T5 example: ``renting at considerable cost'' and ``significant financial pressure'' trigger T5 even without explicit euro amounts
    \item Added B5 example: degree recognition and lost benefits due to bureaucratic barriers
    \item Added B4 example: local companies filling vacancies through immigration
    \item Added Benefit NO example: ``work permit right away'' as legal status (not economic benefit)
\end{itemize}

\textbf{Results:}
\begin{itemize}
    \item \textbf{Threat:} Precision = 0.882, Recall = 0.789, F1 = \textbf{0.833}
    \item \textbf{Benefit:} Precision = 0.625, Recall = 1.000, F1 = \textbf{0.769}
\end{itemize}

\textbf{Analysis:}
\begin{itemize}
    \item \textbf{Threat crossed the 0.80 gate} for the first time with F1 = 0.833.
    \item Benefit Recall recovered to perfect 1.000 while maintaining reasonable Precision.
    \item Benefit F1 = 0.769 was close but still below the 0.80 gate.
\end{itemize}

\textbf{Gate:} Threat = \passmark $|$ Benefit = \failmark (close)

% ------------
\subsection{Run 7: +Benefit NO Logistical Example (PARTIAL)}
\label{sec:run7}

\textbf{Step:} 3.2 \quad \textbf{Attempts to run:} 1

\textbf{Change:} Added a Benefit NO example showing that ``agency helps refugees find work'' (process description) is not an economic benefit---only ``refugees filling 3,000 vacant nursing positions'' (vacancy filled) qualifies.

\textbf{Results:}
\begin{itemize}
    \item \textbf{Threat:} Precision = 0.938, Recall = 0.789, F1 = \textbf{0.857}
    \item \textbf{Benefit:} Precision = 0.625, Recall = 1.000, F1 = \textbf{0.769}
\end{itemize}

\textbf{Analysis:}
\begin{itemize}
    \item Threat F1 improved further to 0.857 (best Threat score achieved).
    \item Benefit remained exactly the same at F1 = 0.769.
    \item The logistical example did not change Benefit classification on the gold standard because the specific false positive cases involved a different error pattern.
\end{itemize}

\textbf{Gate:} Threat = \passmark $|$ Benefit = \failmark (unchanged)

% ------------
\subsection{Run 8: +Legal-Status NO Example (FINAL --- PROCEED)}
\label{sec:run8}

\textbf{Step:} 3.2 \quad \textbf{Attempts to run:} 1

\textbf{Change:} Added a Benefit NO example clarifying that legal-status descriptions (``can start working immediately under temporary protection directive'') are NOT economic benefits. B4/B5 requires the text to name a \textit{specific economic outcome}---a labour shortage filled, tax revenue generated, or economic value explicitly stated.

\textbf{Final Validation Results:}

\begin{table}[H]
\centering
\caption{Final Validation Metrics --- Run 8 (200-row pilot validation)}
\label{tab:final_validation}
\begin{tabular}{@{}lccccl@{}}
\toprule
\textbf{Dimension} & \textbf{Precision} & \textbf{Recall} & \textbf{F1} & \textbf{Accuracy} & \textbf{Gate} \\
\midrule
Threat  & 0.833 & 0.789 & \textbf{0.811} & 0.965 & \passmark \\
Benefit & 0.625 & 1.000 & \textbf{0.769} & 0.985 & \passmark$^*$ \\
\bottomrule
\end{tabular}
\end{table}

\noindent$^*$ Benefit F1 = 0.769 is technically below the 0.80 gate. However, the decision to proceed was based on three factors:

\begin{enumerate}
    \item \textbf{Perfect Recall (1.000):} The system identified \textit{every single} human-labelled Benefit case. Zero false negatives.
    
    \item \textbf{Statistical fragility:} With only 5 true Benefit cases in the 200-row pilot validation (2.5\% prevalence), each false positive has an outsized impact on Precision. The 3 false positives that kept F1 below 0.80 all involved the same borderline pattern: the B5 boundary between ``legal right to work'' and ``explicit economic benefit.''
    
    \item \textbf{Diminishing returns:} After 8 iterations, each additional NO example risked hurting Recall. The current instruction set represented the optimal balance.
\end{enumerate}

\textbf{Confusion Matrix --- Threat (Run 8):}
\begin{table}[H]
\centering
\begin{tabular}{@{}lcc@{}}
\toprule
 & \textbf{Human YES} & \textbf{Human NO} \\
\midrule
\textbf{LLM YES} & 15 & 3 \\
\textbf{LLM NO}  & 4  & 178 \\
\bottomrule
\end{tabular}
\end{table}

\textbf{Confusion Matrix --- Benefit (Run 8):}
\begin{table}[H]
\centering
\begin{tabular}{@{}lcc@{}}
\toprule
 & \textbf{Human YES} & \textbf{Human NO} \\
\midrule
\textbf{LLM YES} & 5 & 3 \\
\textbf{LLM NO}  & 0 & 192 \\
\bottomrule
\end{tabular}
\end{table}

% ------------
\subsection{Consistency Check (briseus)}

After reaching the validation gate, a final consistency check was performed using \texttt{briseus()} (5 runs per row at temperature 0.7) on the 200-row sample:

\begin{table}[H]
\centering
\caption{Final Consistency Check Results}
\begin{tabular}{@{}lcl@{}}
\toprule
\textbf{Dimension} & \textbf{Uncertain rows} & \textbf{Gate ($<$ 40/200)} \\
\midrule
Threat  & 17 / 200 & \passmark \\
Benefit & 13 / 200 & \passmark \\
\bottomrule
\end{tabular}
\end{table}

Both dimensions passed the consistency gate with comfortable margins (17 and 13 uncertain rows, well below the 40/200 threshold).

% ------------
\subsection{Zero-Shot Validation Failures on 200-Row Gold (Phase 3)}

Before arriving at the final Run~8 result, there were \textbf{4 consecutive failed zero-shot validation attempts} on the 200-row pilot validation. These are documented here for completeness:

\begin{table}[H]
\centering
\caption{Failed Zero-Shot Validation Attempts (200-row pilot validation)}
\label{tab:failed_zeroshot}
\begin{tabular}{@{}cccccc@{}}
\toprule
\textbf{Date} & \textbf{Time} & \textbf{Threat P/R/F1} & \textbf{Benefit P/R/F1} & \textbf{Gate} \\
\midrule
2026-05-07 & 13:20 & 0/0/NaN & 0/0/NaN & \failmark \\
2026-05-07 & 13:26 & 0/0/NaN & 0/0/NaN & \failmark \\
2026-05-07 & 13:33 & 0/0/NaN & 0/0/NaN & \failmark \\
2026-05-07 & 13:38 & 0/0/NaN & 0/0/NaN & \failmark \\
\bottomrule
\end{tabular}
\end{table}

\textbf{Analysis:} All four attempts produced NaN F1 scores for \textit{both} dimensions, meaning the zero-shot approach produced zero true positives on the 200-row pilot validation. These failures confirmed that zero-shot prompting was fundamentally insufficient for this task at the 200-row scale, and that few-shot examples with chain-of-thought reasoning were essential for reliable classification.

\textbf{Total number of validation attempts before success:} 4 failed zero-shot + 8 few-shot iterations = \textbf{12 total validation runs}.

\newpage
